\section*{Methods}

\subsection*{Dataset}
We use two datasets: one real benchmark and one synthetic benchmark.

\subsubsection*{Real-data benchmark: colon cancer}
We use the colon cancer dataset, originally introduced by \cite{alon1999}. It is a gene-expression
microarray benchmark built from colon tumor and normal tissue samples.
In the commonly used machine-learning representation, the data matrix contains:
\begin{itemize}[noitemsep]
\item 62 objects (samples),
\item 2,000 gene-expression features,
\item binary labels: \texttt{tumor} (40) and \texttt{normal} (22).
\end{itemize}

The dataset originates from the Alon et al. colon adenocarcinoma study published in 1999 \cite{alon1999}
and has since become a long-standing reference problem for sparse learning and feature selection in
bioinformatics, with representative comparative results in \cite{guyon2002,dudoit2002,li2004}.

\subsubsection*{Synthetic benchmark with known signal strength}
To assess feature recovery directly, we additionally generate a synthetic binary classification dataset with
a known ground-truth definition. The synthetic benchmark is designed to separate predictive performance
from feature-discovery quality.

The guiding idea is simple. The benchmark should contain a small set of truly informative variables, a larger
set of variables that are correlated with the true signal but are not themselves primary signal variables, and
many unrelated noise variables. This lets us ask whether later neurons recover new informative signal or only
move toward redundant and noisy features.

The final synthetic dataset used in this study has:
\begin{itemize}[noitemsep]
\item 100 samples,
\item 2,000 features in total,
\item balanced classes: 50 negative and 50 positive samples,
\item 12 informative features,
\item 36 redundant features,
\item 1,952 noise features.
\end{itemize}

The construction is based on three latent signal sources. For each sample $i$ and source
$b \in \{1,2,3\}$, a latent value $z_{ib}$ is drawn from a standard normal distribution. The three sources have
decreasing strengths
\[
s=(1.2,\;0.9,\;0.6),
\]
so source 1 is strongest and source 3 is weakest. The class score is generated directly from these latent
signals:
\[
r_i=\sum_{b=1}^{3}s_b z_{ib}+\eta_i,
\]
where $\eta_i$ is Gaussian label noise with scale 0.55. The score is thresholded at its median, yielding
exactly 50 negative and 50 positive samples.

Each latent source is then represented in the observed feature matrix by four truly informative features and
twelve correlated redundant features. Thus, the observed signal-related part of the feature matrix has
$3\times4=12$ informative features and $3\times12=36$ redundant features; the remaining 1,952 features are
independent noise. Informative features are indexed first, followed by redundant features, and finally by noise
features.

For an informative feature $j$ assigned to source $b$, the observed value is a noisy measurement of the
corresponding latent source:
\[
x_{ij}=w_j z_{ib}+\epsilon_{ij}.
\]
The feature noise is Gaussian with scale 0.2. The informative weights decrease in absolute value within each
signal source and alternate in sign. Thus, informative features have a known order of importance both across
signal sources and within each source, while the benchmark does not reduce to selecting only positively
correlated variables. The exact informative-feature weights and feature-index mapping are fixed by the
benchmark definition.

Redundant features are also assigned to one signal source. Each redundant feature is generated as a correlated
surrogate of one informative feature from that source and also receives a smaller contribution from the same
latent value $z_{ib}$. The redundant-feature noise scale is 0.3. These variables can be
predictive because they are correlated with the signal, but they are not counted as truly informative in the
ground truth.

Class labels are therefore generated from the three latent signal sources, not from all 2,000 observed features. This
preserves a known latent signal structure while allowing the observed feature matrix to contain informative,
redundant, and purely noisy variables.

After generation, all feature columns are standardized by z-scoring. As a result, feature variances are on a
common scale, but the effective signal-to-noise ratios remain different across informative, redundant, and
noise features.

The benchmark definition specifies the exact indices of informative, redundant, and noise features. This makes it
possible to evaluate not only predictive accuracy, but also how many of the 12 ground-truth informative
features are recovered by each neuron and by the cumulative sequence of neurons.

\subsection*{Model family}
The core model is a complex-layer classifier that sequentially trains sparse linear neurons and removes selected
features after each neuron. This construction follows the Complex Layers of Formal Neurons line of work
\cite{bobrowski2021,bobrowski2022}. For analysis, we treat neuron depth $n$ as a model dimension in its own
right. For each depth, we inspect both the prediction of the $n$-th neuron itself and the cumulative prediction
obtained by majority voting over the first $n$ neurons.

The label \texttt{cpl} refers to the convex piecewise-linear base learner used in this study. Its
objective combines hinge loss with L1 regularization; both terms are convex piecewise-linear, so the resulting
sparse linear neuron can be trained as a linear programming problem. It is evaluated inside the same
complex-layer protocol as two standard sparse linear baselines:
\begin{itemize}[noitemsep]
\item \texttt{cpl}: convex piecewise-linear neuron trained with an LP-based hinge + L1 formulation.
\item \texttt{svm}: LinearSVC (squared hinge loss + L1 regularization).
\item \texttt{logreg}: logistic regression with L1 regularization.
\end{itemize}
The use of L1-type regularization is motivated by its long-standing role in sparse estimation and feature
selection for high-dimensional problems \cite{tibshirani1996lasso,ng2004l1l2}. For logistic regression, the
computational perspective is closely related to regularization-path methods such as \texttt{glmnet}
\cite{friedman2010glmnet}, while the practical large-scale linear classification setting is represented by
\texttt{LIBLINEAR} \cite{fan2008liblinear}.

\subsection*{Evaluation protocol}
For both benchmarks, we use repeated stratified cross-validation with a fixed random seed.
Both benchmarks reported in this manuscript use 10 folds and 10 repeats, giving 100 train/test splits for
each dataset and model configuration.
The split-wise protocol is important in this setting because feature selection bias is a well-known failure
mode in high-dimensional gene-expression studies when selection and evaluation are not strictly separated
\cite{ambroise2002,krawczuk2016feature}.

\subsection*{Computational environment}
All experiments were implemented in Python and are distributed with the manuscript repository. The submitted
artifact specifies Python 3.13 or 3.14 with Poetry-managed dependencies; the final local verification used
Python 3.14.3 and Poetry 2.3.2. Package versions are pinned in \texttt{poetry.lock}, including
\texttt{numpy}, \texttt{scipy}, \texttt{scikit-learn}, \texttt{pandas}, and \texttt{matplotlib}. The paper-scale
runs were executed on a MacBook Air (Apple M2, 2022; 16 GB RAM) running macOS Tahoe 26.5.1, and each run
records wall-clock time and peak resident memory in its \texttt{run\_stats.json} metadata. The repository
scripts are the authoritative workflow for regenerating raw predictions, metric tables, and figures from the
input datasets.

\subsection*{Regularization calibration}
Because the three base learners use different optimization schemes and induce sparsity at different rates,
their regularization parameters are not directly comparable on a shared numerical scale. To make the comparison
more interpretable, we calibrate two representative regularization settings for each model on each dataset.

The calibration procedure is intentionally simple. For a given dataset, model, and value of $C$, we fit only
the first neuron and record the number of selected features. We then choose two values of $C$:
\begin{itemize}[noitemsep]
\item a \textbf{low-feature} setting targeting approximately five selected features in the first neuron;
\item a \textbf{high-feature} setting targeting approximately twenty-five selected features in the first neuron.
\end{itemize}

This calibration does not make the models equivalent, and it is not intended to support a strict ranking of
the base learners. Its purpose is to place each learner in two interpretable feature-count regimes and then
study how that learner behaves as successive neurons are added. Realized first-neuron widths can still differ
across learners under cross-validation, especially when the sparsity path changes abruptly with $C$.
This choice follows the broader methodological view that comparisons between sparse learners should distinguish
the induced first-neuron width from the raw numerical value of the regularization parameter itself
\cite{ng2004l1l2,saeys2007review}.

\begin{table}[ht]
\centering
\begin{tabular}{llrrrr}
\toprule
Dataset & Model & $C_{\mathrm{low}}$ & Features & $C_{\mathrm{high}}$ & Features \\
\midrule
Colon cancer & \texttt{cpl} & 0.03 & 6 & 0.08 & 24 \\
Colon cancer & \texttt{svm} & 1.2 & 5 & 6.0 & 25 \\
Colon cancer & \texttt{logreg} & 3.5 & 6 & 6.4 & 26 \\
Synthetic benchmark & \texttt{cpl} & 0.035 & 4 & 0.05 & 27 \\
Synthetic benchmark & \texttt{svm} & 2.0 & 6 & 3.5 & 26 \\
Synthetic benchmark & \texttt{logreg} & 5.0 & 7 & 9.2 & 25 \\
\bottomrule
\end{tabular}
\caption{Calibration of regularization strength by first-neuron feature count. For each dataset and model, the
parameter $C_{\mathrm{low}}$ was selected to produce roughly five selected features in the first neuron, and
$C_{\mathrm{high}}$ to produce roughly twenty-five selected features. The reported feature counts come from a
single fit of the first neuron on the full dataset and are used only to define comparable low-feature and
high-feature regimes before repeated cross-validation.}
\label{tab:regularization_calibration}
\end{table}

For each split:
\begin{enumerate}[noitemsep]
\item The model is trained only on the training fold.
\item A sequence of up to 31 neurons is fitted.
\item After each neuron fit, features with non-zero coefficients are removed from the candidate pool
for subsequent neurons.
\item Predictions are computed on the held-out test fold.
\end{enumerate}

For each split and neuron depth $n$:
\begin{itemize}[noitemsep]
\item \textbf{vote prediction}: majority vote over the first $n$ neurons;
\item \textbf{weighted vote prediction}: vote over the first $n$ neurons with weights estimated only from the
training fold;
\item \textbf{neuron prediction}: prediction of the $n$-th neuron only.
\end{itemize}

For the weighted vote, each neuron receives a non-negative quality score derived from its training-fold
balanced accuracy:
\[
q_j=\max\left(0,\mathrm{BA}^{\mathrm{train}}_j-0.5\right),
\qquad
w_j=\frac{q_j}{\sum_{\ell=1}^{n} q_\ell}.
\]
If all $q_j$ are zero, uniform weights are used. The chance-level offset of 0.5 makes neurons with no
above-chance balanced accuracy contribute zero weight. Because all weights are computed on the training fold,
this procedure does not use held-out labels.

Then, for each neuron index $n$, metrics are averaged over all 100 splits.

For each neuron depth, we report the mean accuracy of the majority vote, the mean held-out accuracy of the
training-quality weighted vote, the accuracy of the $n$-th neuron alone, and the mean number of features
selected by that neuron.

For the synthetic benchmark, the known set of 12 informative features also allows feature-recovery analysis.
Feature precision measures what fraction of the selected features are truly informative, while feature recall
measures what fraction of the 12 informative features have been recovered. We use both per-neuron and
cumulative recall to assess whether later neurons add new signal or mainly move toward redundant and noisy
variables.
This emphasis goes beyond pure prediction accuracy and is consistent with prior arguments that feature
selection quality and stability deserve explicit evaluation in high-dimensional settings
\cite{saeys2007review,nogueira2018stability,lukaszuk2024importance}.

For binary labels $y_i \in \{-1,+1\}$ and linear score $f(x_i)=w^\top x_i + b$, the objective forms are:
\[
\texttt{cpl}: \quad
\min_{w,b}\ \frac{1}{N}\sum_{i=1}^{N}\max\!\left(0,\,1-y_i f(x_i)\right) + \lambda\lVert w\rVert_1
\]
\[
\texttt{svm}: \quad
\min_{w,b}\ \frac{1}{N}\sum_{i=1}^{N}\max\!\left(0,\,1-y_i f(x_i)\right)^2 + \lambda\lVert w\rVert_1
\]
\[
\texttt{logreg}: \quad
\min_{w,b}\ \frac{1}{N}\sum_{i=1}^{N}\log\!\left(1+\exp(-y_i f(x_i))\right) + \lambda\lVert w\rVert_1
\]

All metrics are reported as mean and standard deviation across splits.

The nominal maximum number of neurons is set to 31; training may stop earlier if no additional non-zero
feature set can be produced.
